\section{Conclusion}

This paper is the first, to our knowledge, to report empirical accounts of how laboratory and field scientific practitioners in embodied and improvisational contexts envision the role of AI. Through on-site interviews and speculative design sessions at 12 scientific practitioners' workplaces, we surfaced three barriers to adoption: experimental setups being too high-stakes to risk AI errors, constrained physical environments, and AI's lack of embodied tacit knowledge. Thus, today AI is mostly used for desk-based science work, but our participants were able to sketch ideas for future tools that scaffold memory and attention during hands-on practice and extend across modalities from robots to lab-wide systems. These speculative designs highlight the importance of tacit knowledge~\cite{polanyi2009tacit}, physical constraints, and spatially distributed workflows. In sum, we aim to shift the focus of AI for science from automating reasoning and computational work at the desk to infrastructure that sustains the embodied physical conditions upon which creative reasoning and discovery depend.
